\subsection{An analytic HS supporting inequality}\label{sec:fqe-hs-polynomial}

At the upper anchor $q_u=8/5$, choosing $\beta=1/12$ makes the
fixed-purity optimisation a bound on $H_{8/5}-t$. We first show that its
maximiser has one large eigenvalue, then prove the resulting one-variable
supporting inequality. The outer dual is nonnegative for $F\le3/20$,
which determines the split between the two anchor regions.
Throughout this subsection put $q=8/5$, $r=q-1=3/5$ and $L=\ln2$.

\begin{lemma}[A third-moment tilt at fixed purity; proved]\label{lem:fqe-hs-tilt}
Among four-point probability vectors of a fixed purity $p\le9/16$,
$H_{8/5}(\lambda)-\sum_i\lambda_i^3$ is maximised only by the one-large spectrum
\[
 \lambda(s)=\bigl((1+3s)/4,(1-s)/4,(1-s)/4,(1-s)/4\bigr),
 \qquad p=(1+3s^2)/4,
\]
up to permutation (uniform when $p=1/4$).
\end{lemma}
\begin{proof}
Let $Z=\sum_i\lambda_i^q$ and let $a,b,c$ be three distinct entries.
Their mean is at most $1/3$, so the divided-difference bound~\eqref{eq:fqe-dd-power} gives
$-[a,b,c]x^r\ge\frac12r(1-r)3^{2-r}$. Also
$Z=\sum_i\lambda_i\lambda_i^r\le p^r$ by concavity of $x^r$.
Consequently
\[
 \eta(a,b,c):=-\frac{q[a,b,c]x^r}{3rLZ}
 \ge\frac8{75L}\frac{3^{7/5}}{(9/16)^{3/5}}
 >\frac{16}{105}(3\cdot2^{12})^{1/5}>1.
 \tag{HTS}\label{eq:fqe-hs-tilt-slope}
\]
The last comparisons use $L<7/10$ and
$105^5=12762815625<3\cdot2^{32}=12884901888$.

Lemma~\ref{lem:fqe-low-tilt-reduction}, with $\omega=1$, reduces the maximum
to the pattern $(a,b,c,c)$. Write it as
\[
 (1/4+u+d,1/4+u-d,1/4-u,1/4-u),\qquad
 x=p-1/4=4u^2+2d^2,\quad 0\le d\le2u.
\]
Decrease $u$ to $u_0=\sqrt{x/12}$ at fixed $x$, setting
$d=\sqrt{(x-4u^2)/2}$. Ordering and positivity persist, and $b=c$ at
the endpoint: it is the one-large spectrum.
Along this path
\[
 t=\frac1{16}+\frac34x+3xu-12u^3,\qquad
 \frac{dt}{du}=3x-36u^2\le0.
\]
As in the preceding chart reduction, $dH_q/dt=\eta(a,b,c)>1$.
Thus decreasing $u$ strictly increases $H_q-t$ unless the spectrum
is already one-large. Continuity handles zero or coalescing entries.
At equality the strict paths force every determinant-maximal replacement to be
one-large. That replacement preserves $p,t$, so equality in the sharp third-moment bound
(Lemma~\ref{lem:fqe-third-moment}) forces the original spectrum to be the same.
The uniform case $p=1/4$ is immediate.
\end{proof}

\begin{lemma}[HS supporting inequality; proved]\label{lem:fqe-hs-polynomial}
At $q=q_u=8/5$ put $u=3^{3/5}$ and
\[
 k=\frac{16}{3L}\frac{u-1}{u+1},\qquad
 \beta=\frac1{12},\qquad \gamma=\frac{10}3,\qquad
 \alpha=\frac32k+\frac{65}{36}.
\]
Every four-point probability vector with $p\le9/16$ satisfies
$\Lambda_{(\alpha,\beta,\gamma)}(\lambda)\le h_{8/5}$,
with equality only at HS. For $0\le\varepsilon\le1/9$ and
$0\le F\le3/20$,
\[
 \ell_{(\alpha,\beta,\gamma)}(\varepsilon,F)
 \ge\frac{149}{324}\varepsilon
       +F\left(\frac1{12}-\frac59F\right)\ge0.
\]
\end{lemma}
\begin{proof}
At fixed $p$, the only variable part of $\Lambda$ is $H_q-t$.
Lemma~\ref{lem:fqe-hs-tilt} reduces its maximum to $\lambda(s)$.
Write $f(p)=H_q(\lambda(s))$, $p_* =1/3$, and
\[
 A(s)=\frac92(s-\tfrac13)^2(s+\tfrac16),\qquad
 W(p)=f(p)-\beta A(s)-\gamma(p-p_*)^2.
\]
Direct moment evaluation gives
\[
 \Lambda(\lambda(s))=W(p)+\frac32k(p-p_*),\qquad
 W(p_*)=h_q,\qquad W'(p_*)=-\frac32k,
 \qquad \frac{d^2 A}{dp^2}=\frac6s.
\]
We show that $W$ is first convex and then concave, allowing either
interval to be empty. The needed scalar calculation is explicit below.
Put
\[
 x=\left(\frac{1+3s}{1-s}\right)^{1/5},\quad
 J=\frac{L f''(p)}{2\gamma+6\beta/s},\quad
 h_4(x)=1+x+x^2+x^3+x^4.
\]
Since $p\le9/16$, $s<2/3$, so $1<x<2$ and $0<x^5-1<8$.
Differentiating
\[
 L\frac{dH_q(\lambda(s))}{ds}
 =-\frac{2(x^5+3)(x^3-1)}{x^8+3}
\]
gives the identity
\[
 \frac{dJ}{dx}
 =-\frac{8(x^5+3)^3 Q(x)}
 {5x^3(43x^5-31)^2(x^8+3)^3h_4(x)^3}.
 \tag{HCJ}\label{eq:fqe-hs-curvature-identity}
\]
Here, with $y=x-1$ and $R=x^5-1$,
$Q(x)=\sum_{j=0}^7R^jP_j(y)$, where the five columns below are the
coefficients of $1,y,y^2,y^3,y^4$ in $P_j$.
\[
\begin{array}{c|rrrrr}
 j&1&y&y^2&y^3&y^4\\\hline
% HD-CURVATURE-BEGIN
0&225600&457600&471040&240000&48640\\
1&582160&1108000&1118304&557584&110864\\
2&690608&1195520&1018512&418768&68128\\
3&756896&1304520&1151162&522393&96703\\
4&516296&613195&462027&220948&46208\\
5&181031&120245&477&-23717&-6797\\
6&-259&2795&1978&774&129\\
7&43&0&0&0&0
% HD-CURVATURE-END
\end{array}
\]
This displayed identity proves the sign without a parameter certificate:
$P_0,\ldots,P_4,P_7$ have positive coefficients, and, for $0\le y<1$,
\[
 P_5(y)\ge181031-23717-6797=150517,\qquad P_6(y)\ge-259.
\]
Thus $P_5+RP_6\ge150517-8\cdot259>0$. Hence $Q>0$ and $J$ strictly decreases.
As $W''=(2\gamma+6\beta/s)(J/L-1)$, the asserted curvature pattern follows.

To turn this curvature pattern into a supporting line, subtract the
tangent at $p_*$ and write
$D(p)=W(p)-W(p_*)-W'(p_*)(p-p_*)$.
Then $D(p_*)=D'(p_*)=0$, and $D'$ first increases and then decreases.
If $D(1/4)<0$, the contact $p_*$ must lie in the strictly concave portion:
otherwise $D'<0$ throughout its left side, which would give
$D(1/4)>0$. In the concave portion, the tangent strictly
dominates to the right of $p_*$; to its left, $D$ has no maximum except
at $1/4$ or $p_*$. Thus $D<0$ away from $p_*$.
In the present case the
left-endpoint gap is strictly positive, so $D(1/4)<0$:
\[
 \Delta=-D(1/4)=h_q+\frac k8-2+\frac\beta{12}+\frac\gamma{144}
 =h_q+\frac k8-2+\frac{13}{432}>\frac1{5000}.
 \tag{HUC}\label{eq:fqe-hs-uniform-contact}
\]
For completeness, this small constant has an elementary rational proof.
Let $u_0=19331/10000$, $v=1/(2u_0+1)$ and
\[
 L_-=2\sum_{j=0}^5\frac{1}{(2j+1)3^{2j+1}},\qquad
 B_+=2\sum_{j=0}^4\frac{v^{2j+1}}{2j+1}
             +\frac{2v^{11}}{11(1-v^2)}.
\]
The formula $h_q=8/3-(5/(3L))\ln(1+1/u)$ and the atanh series give $L>L_-$ and
$\ln(1+1/u_0)<B_+$. Since $u_0^5<27$, $u>u_0$; the function
$-5\ln(1+1/u)+2(u-1)/(u+1)$ increases with $u>0$.
The purely rational comparison
\[
 \frac{301}{144}L_- -5B_+ +\frac{2(u_0-1)}{u_0+1}
 >\frac{21}{50000}
\]
therefore gives $3L\Delta>21/50000>3L/5000$, using $L<7/10$.
This proves~\eqref{eq:fqe-hs-uniform-contact}.
The tangent bound is strict away from $p_*$ by the strict curvature
calculation. Thus $\Lambda\le h_q$, with equality only at HS.

Finally $u>19/10$ and $L<7/10$ give
$k>480/203>7/3$. Substitution into the outer dual yields
\[
 \ell=\left(\alpha-\gamma(\tfrac{17}{12}+\tfrac\varepsilon3)\right)\varepsilon
          +F\left(\tfrac1{12}-\tfrac59F\right)
 \ge\frac{149}{324}\varepsilon+F\left(\tfrac1{12}-\tfrac59F\right).
\]
The last term is nonnegative for $F\le3/20$, as required.
All global maximisation and sign arguments above are analytic;
HD checks the displayed identities and rational comparisons.
\end{proof}
